\documentclass[11pt]{article}
\usepackage{amsmath, amssymb}
\usepackage[margin=1in]{geometry}
\usepackage{fancyhdr}
\usepackage{hyperref}

\hypersetup{
    colorlinks=true,
    linkcolor=blue,
    urlcolor=blue,
    citecolor=green
}

% ---- Fill these in ----
\newcommand{\courseName}{NEU 314}
\newcommand{\taName}{Keith Murray}
\newcommand{\sectionInfo}{Precept B01}
% ------------------------

\title{Precept: Bayesian Exercises}
\date{\today}

\pagestyle{fancy}
\fancyhf{}
\lhead{\courseName}
\chead{\sectionInfo}
\rhead{\today}
\cfoot{\thepage}

\begin{document}

\maketitle
\thispagestyle{fancy}

\section*{Problems}
Remember from earlier today Bayes' rule
\[
P(x\mid y) = \frac{P(y\mid x)P(x)}{P(y)}
\]

\begin{enumerate}
    \item \textbf{The positive test.} A rare condition affects $1\%$ of the population. A test for it comes back positive $99\%$ of the time when someone has the condition, and gives a false positive $5\%$ of the time when someone does not.
    \begin{enumerate}
        \item You take the test and it comes back positive. What is the probability that you actually have the condition?
    \end{enumerate}

    \item \textbf{The trick coin.} A bag holds two coins. One is fair, and the other has heads on both sides. You pull a coin out at random and flip it.
    \begin{enumerate}
        \item It lands heads. What is the probability that you pulled out the two-headed coin?
        \item You flip the same coin again and it lands heads again. What is the probability now?
    \end{enumerate}

    \item \textbf{The two-child puzzle.} A family has two children. Assume each child is equally likely to be a girl or a boy, independently of the other.
    \begin{enumerate}
        \item You learn that \emph{at least one} of the children is a girl. What is the probability that both are girls?
        \item Instead, you learn that the \emph{older} child is a girl. What is the probability that both are girls?
    \end{enumerate}
\end{enumerate}

\newpage
\section*{Solution}

\subsection*{Question 1}
Let $C$ mean ``has the condition'' and $\bar{C}$ mean ``does not.'' We are given
\[
P(C)=0.01,\qquad P(+\mid C)=0.99,\qquad P(+\mid \bar{C})=0.05
\]

\subsubsection*{Part (a)}
We can just use Bayes' rule!
\[
P(D\mid +)=\frac{P(+\mid D)\,P(D)}{P(+)}
\]
By substituting in the values given to us, we can compute $P(+)$ as 
\[
P(+)=P(+,D)+P(+,\bar{D})=P(+\mid D)P(D)+P(+\mid\bar{D})P(\bar{D})=(0.99)(0.01)+(0.05)(0.99)=0.0594
\]
Therefore, we can write
\[
P(D\mid +)=\frac{0.0099}{0.0594}=\frac{1}{6}\approx 0.17
\]
Funny how that's not 99\%, right?

\subsection*{Question 2}
Let $T$ be ``two-headed coin'' and $F$ be ``fair coin,'' so $P(T)=P(F)=\tfrac{1}{2}$, with $P(H\mid T)=1$ and $P(H\mid F)=\tfrac{1}{2}$.

\subsubsection*{Part (a)}
This is just Bayes' rule again!
\[
P(T\mid H)=\frac{P(H\mid T)P(T)}{P(H)}
\]
To compute $P(H)$, we can use marginalization and the product rule to get
\[
P(H)=P(H\mid T)P(T)+P(H\mid F)P(F)=(1)\left(\tfrac{1}{2}\right)+\left(\tfrac{1}{2}\right)\left(\tfrac{1}{2}\right)=\tfrac{3}{4}
\]
so that
\[
P(T\mid H)=\frac{P(H\mid T)P(T)}{P(H)}=\frac{1/2}{3/4}=\frac{2}{3}
\]

\subsubsection*{Part (b)}
Given the coin, the flips are independent, so $P(HH\mid T)=1$ and $P(HH\mid F)=\tfrac{1}{4}$. Then
\[
P(T\mid HH)=\frac{(1)\left(\tfrac{1}{2}\right)}{(1)\left(\tfrac{1}{2}\right)+\left(\tfrac{1}{4}\right)\left(\tfrac{1}{2}\right)}=\frac{1/2}{5/8}=\frac{4}{5}
\]
Each head makes us more confident in the two-headed coin. Notice that a single tails would immediately drop $P(T)$ to $0$.

\subsection*{Question 3}
Listing the children as (older, younger), the four outcomes $GG$, $GB$, $BG$, $BB$ each have probability $\tfrac{1}{4}$.

\subsubsection*{Part (a)}
Again! This is Bayes' rule!
\[
P(GG\mid \text{at least one } G)=\frac{P(\text{at least one } G\mid GG)\cdot P(GG)}{P(\text{at least one } G)}
\]
Given that there is ``at least one girl'', we know that there are three possible states of the world $\{GG,GB,BG\}$, so 
\[
P(\text{at least one } G)=\frac{3}{4}
\]
Notice that $P(\text{at least one } G\mid GG)=1$ and $P(GG)=\frac{1}{4}$. Hence, we get
Since $GG$ is contained in this event, $P(GG,\text{at least one } G)=P(GG)=\tfrac{1}{4}$. By the product rule,
\[
P(GG\mid \text{at least one } G)=\frac{P(\text{at least one } G\mid GG)\cdot P(GG)}{P(\text{at least one } G)}=\frac{1/4}{3/4}=\frac{1}{3}
\]

\subsubsection*{Part (b)}
We can intuitively reason about this one. Since the ``older child is a girl'', there are two possible states of the world $\{GG,GB\}$. Given they are both equal probability, the probability that both children are girls is $\frac{1}{2}$.

\end{document}